\documentclass[11pt,a4paper]{article}
\usepackage[margin=19mm,top=17mm,bottom=17mm]{geometry}
\usepackage{amsmath,amssymb,amsfonts}
\usepackage{mathptmx}
\usepackage{enumitem}
\usepackage{fancyhdr}
\usepackage{array,tabularx,booktabs}
\usepackage[hidelinks]{hyperref}
\setlength{\parindent}{0pt}\setlength{\parskip}{4pt}
\setlist[enumerate,1]{leftmargin=1.1em,itemsep=3pt,parsep=2pt}
\setlist[itemize]{leftmargin=1.4em,itemsep=1pt}
\pagestyle{fancy}\fancyhf{}
\fancyhead[C]{\footnotesize \textbf{Question Paper}}
\fancyfoot[C]{\footnotesize Page \thepage}
\renewcommand{\headrulewidth}{0.4pt}
\begin{document}
\vspace{2mm}
\noindent \textbf{Marks : 28} \hfill \textbf{Time : 2 : 30 hours}
\noindent\rule{\linewidth}{1.1pt}\vspace{-1mm}\noindent\rule{\linewidth}{0.5pt}
{\centering \Large \textbf{DEFINITE INTEGRALS PRACTICE PAPER}\par}
{\centering JEE $|$ Mathematics\par}
\noindent\rule{\linewidth}{1.1pt}\vspace{-1mm}\noindent\rule{\linewidth}{0.5pt}
\vspace{1mm}
\noindent \textbf{Duration:} Not specified \quad \textbf{Total Marks:} 28 \quad \textbf{Questions:} 10\par
\vspace{2mm}

\begin{tabularx}{\linewidth}{|X|X|}\hline
Candidate Name: \rule{6cm}{0.4pt} & Roll Number: \rule{3.5cm}{0.4pt} \\\hline
\end{tabularx}

\textbf{Instructions}
\begin{itemize}[topsep=2pt]
\item Read every question carefully before selecting an answer.
\item For single-correct questions, mark one option only.
\item Use the answer sheet or space provided by the invigilator.
\end{itemize}
\begin{enumerate}[resume]
\item Let \(k\in\mathbb{R}\) and define
\[
J(k)=\int_0^1\left[\frac{x^2}{1+x^2}+k\frac{(1-x)^2}{1+(1-x)^2}\right]dx.
\]
Suppose that \(J(k)\) has the same value as the integral obtained by replacing \(x\) with \(1-x\) throughout the integrand. Then \(k\) can be: \hfill [4 marks]
{\renewcommand{\arraystretch}{1.8}
\begin{tabularx}{\linewidth}{X X}
(A) \quad \(0\) only & (B) \quad \(1\) only \\[10pt]
(C) \quad \(-1\) only & (D) \quad Every real number \\[10pt]
\end{tabularx}}
\item Let $f$ be an integrable function on $[0,2]$ satisfying $f(x)+f(2-x)=x(2-x)+1$ for every $x\in[0,2]$. If
\[
J=\int_0^2 x(2-x)f(x)\,dx,
\]
then the exact value of $J$ is: \hfill [4 marks]
{\renewcommand{\arraystretch}{1.8}
\begin{tabularx}{\linewidth}{X X X X}
(A) \quad $\frac{6}{5}$ & (B) \quad $\frac{12}{5}$ & (C) \quad $\frac{3}{5}$ & (D) \quad $\frac{8}{5}$ \\[10pt]
\end{tabularx}}
\item Let $f$ be continuous on $[0,1]$ and satisfy
\[
f(x)+f(1-x)=1+x(1-x),\qquad 0\le x\le 1.
\]
If
\[
J=\int_0^1 x(1-x)f(x)\,dx,
\]
then the value of $J$ is: \hfill [4 marks]
{\renewcommand{\arraystretch}{1.8}
\begin{tabularx}{\linewidth}{X X X X}
(A) \quad $\frac{1}{12}$ & (B) \quad $\frac{1}{10}$ & (C) \quad $\frac{1}{8}$ & (D) \quad $\frac{1}{6}$ \\[10pt]
\end{tabularx}}
\item Let \[I=\int_{0}^{1}\frac{(1+x)^7}{(1+x)^7+(2-x)^7}\,dx.\] Then the value of $I$ is: \hfill [4 marks]
{\renewcommand{\arraystretch}{1.8}
\begin{tabularx}{\linewidth}{X X X X}
(A) \quad $\frac{1}{4}$ & (B) \quad $\frac{1}{2}$ & (C) \quad $\frac{7}{8}$ & (D) \quad $1$ \\[10pt]
\end{tabularx}}
\item Let \(\rho:[0,L]\to[0,\infty)\) be a continuous function, not identically zero, satisfying \(\rho(x)=\rho(L-x)\) for every \(x\in[0,L]\). If the center of mass of the corresponding one-dimensional distribution is \[\bar{x}=\frac{\int_0^L x\rho(x)\,dx}{\int_0^L \rho(x)\,dx},\] then \(\bar{x}\) is equal to: \hfill [1 marks]
{\renewcommand{\arraystretch}{1.8}
\begin{tabularx}{\linewidth}{X X}
(A) \quad \(\dfrac{L}{4}\) & (B) \quad \(\dfrac{L}{2}\) \\[10pt]
(C) \quad \(\dfrac{2L}{3}\) & (D) \quad It depends on the particular form of \(\rho\) \\[10pt]
\end{tabularx}}
\item Let \(\lambda\in\mathbb{R}\) and define
\[
I(\lambda)=\int_{0}^{1}\frac{x^2+\lambda x}{1+x^2}\,dx.
\]
The value(s) of \(\lambda\) for which
\[
I(\lambda)=\int_{0}^{1}\frac{(1-x)^2+\lambda(1-x)}{1+(1-x)^2}\,dx
\]
hold are: \hfill [1 marks]
{\renewcommand{\arraystretch}{1.8}
\begin{tabularx}{\linewidth}{X X}
(A) \quad \(\lambda=0\) only & (B) \quad \(\lambda=1\) only \\[10pt]
(C) \quad \(\lambda=-1\) only & (D) \quad Every real value of \(\lambda\) \\[10pt]
\end{tabularx}}
\item Let \(\lambda\in\mathbb{R}\) satisfy
\[
\int_0^1 (x-\lambda)\sqrt{x(1-x)}\,[1+x(1-x)]\,dx=0.
\]
Then the value of \(\lambda\) is: \hfill [4 marks]
{\renewcommand{\arraystretch}{1.8}
\begin{tabularx}{\linewidth}{X X X X}
(A) \quad \(0\) & (B) \quad \(\frac14\) & (C) \quad \(\frac12\) & (D) \quad \(1\) \\[10pt]
\end{tabularx}}
\item Let $g:[0,1]\to\mathbb{R}$ be defined by $g(x)=1+x(1-x)$. If $c\in(0,1)$ satisfies
\[
\int_0^c g(x)\,dx=\int_c^1 g(x)\,dx,
\]
then the value of $c$ is: \hfill [1 marks]
{\renewcommand{\arraystretch}{1.8}
\begin{tabularx}{\linewidth}{X X X X}
(A) \quad $\frac{1}{3}$ & (B) \quad $\frac{1}{2}$ & (C) \quad $\frac{2}{3}$ & (D) \quad $\frac{3}{4}$ \\[10pt]
\end{tabularx}}
\item Let $f$ be any integrable function on $[1,5]$. If the relation
\[
\int_{1}^{5}(x-c)f(x)\,dx=-\int_{1}^{5}(x-c)f(6-x)\,dx
\]
holds for every such function $f$, then the value of the real constant $c$ is: \hfill [1 marks]
{\renewcommand{\arraystretch}{1.8}
\begin{tabularx}{\linewidth}{X X X X}
(A) \quad $1$ & (B) \quad $2$ & (C) \quad $3$ & (D) \quad $4$ \\[10pt]
\end{tabularx}}
\item Let $a$ be a positive real number satisfying
\[
\int_0^1 \frac{x^a}{x^a+(1-x)^a}\,dx=\frac{a}{4}.
\]
The value of $a$ is: \hfill [4 marks]
{\renewcommand{\arraystretch}{1.8}
\begin{tabularx}{\linewidth}{X X X X}
(A) \quad $\frac{1}{2}$ & (B) \quad $1$ & (C) \quad $2$ & (D) \quad $4$ \\[10pt]
\end{tabularx}}
\end{enumerate}
\end{document}